% !TeX program = pdflatex
% papers/corpo_universal.tex — o Corpo Universal: as cinco operações, o
% observador, a transformada e a membrana de Dirac. O Corpo de Peano é
% instância. NENHUMA afirmação sem medidor: cada teorema cita a bateria.
% Auto-contido: pdflatex corpo_universal.tex
\documentclass[11pt,a4paper]{article}
\usepackage[utf8]{inputenc}
\usepackage[T1]{fontenc}
\usepackage[portuguese]{babel}
\usepackage{amsmath,amssymb,amsthm}
\usepackage[margin=2.5cm]{geometry}
\usepackage{booktabs}
\usepackage{xcolor}
\usepackage[hidelinks]{hyperref}
\newtheorem{definicao}{Definição}
\newtheorem{teorema}{Teorema}
\newtheorem{proposicao}{Proposição}
\newtheorem{corolario}{Corolário}
\newtheorem{observacao}{Observação}
\newcommand{\code}[1]{\texttt{#1}}
% ─────────────────────────────────────────────────────────────────────────────
%  papers/gkcapa.tex — a capa do Reino, mínima, para os papers em `article`.
%
%  O estilo.tex completo é do `report` (títulos de capítulo, skak, …). Aqui fica
%  só o que a capa precisa, para o paper continuar a compilar sozinho com
%  pdflatex. O tradutor próprio (tests/tex) carrega o estilo.tex da raiz / o
%  symlink papers/estilo.tex e avalia o mesmo `\gkcapa`.
% ─────────────────────────────────────────────────────────────────────────────
\definecolor{ouro}{HTML}{B8912F}
\definecolor{ouroclaro}{HTML}{F3E9CE}
\definecolor{tinta}{HTML}{1E1B16}
\definecolor{regua}{HTML}{6E6858}
\providecommand{\gknota}{\fontsize{7.62}{11.05}\selectfont}
\providecommand{\gktexto}{\fontsize{10.50}{15.22}\selectfont}
\providecommand{\gksub}{\fontsize{12.33}{17.87}\selectfont}
\providecommand{\gksec}{\fontsize{14.47}{20.98}\selectfont}
\providecommand{\gkcap}{\fontsize{16.99}{24.63}\selectfont}
\providecommand{\gktit}{\fontsize{23.42}{33.95}\selectfont}
\providecommand{\Prj}{\mathbb{P}}
\providecommand{\Trf}{\mathcal{T}}
\providecommand{\gkcapa}[3]{%
\title{\vspace{-0.6cm}
{\color{ouro}\rule{\textwidth}{1.4pt}}\\[6mm]
\textbf{\gktit\color{tinta} Reino Dourado}\\[3mm]{\gksec\itshape\color{regua} Golden Kingdom}\\[8mm]
{\gkcap\scshape\color{ouro} #1}\\[5mm]
{\gksub\color{regua} #2}\\[2mm]
{\gktexto\color{regua}\itshape #3}\\[7mm]
{\color{ouro}\rule{\textwidth}{0.6pt}}\\[9mm]{\gknota\color{regua}\begin{minipage}{\textwidth}\setlength{\parindent}{0pt}%
\textbf{Aviso de Isenção Algorítmica:} O universo, as leis e as entidades contidas nestas páginas ---
incluindo felinos matriciais, aves de rapina algébricas e amuletos de controle de fluxo --- são
construções de pura ficção formal. Esta obra não descreve a realidade física, não serve como
engenharia reversa do cosmos e não constitui doutrina religiosa. Qualquer semelhança com bugs reais do seu
sistema operacional é mera coincidência (ou falta de reza). Prossiga por sua própria conta e
risco.\end{minipage}}}
\author{}\date{}}

\begin{document}
\gkcapa{O Corpo Universal}
       {as cinco operações, o observador, a membrana --- o Peano é instância}
       {tudo aqui tem medidor; a visão unificada constrói-se por elos, e cada elo aguenta peso}
\maketitle

\begin{abstract}
O \textbf{Corpo Universal} é o quadro dos operadores que a bateria já
mediu, organizado pela cadeia
\[
\boxed{
\text{corpo}\to\text{norma}\to\text{ordem}\to\text{observador}
\to\text{transformação}\to\text{dual}.
}
\]
As \textbf{cinco operações} --- Soma, Multiplicação, Divisão, Dual,
Inversão --- entram cada uma com a sua conservação \emph{medida}; o
\textbf{observador} é a escada $I_1\prec I_2\prec I_3\prec\cdots$; a
\textbf{transformada universal} é a avaliação nas folhas, realizada
inteira pela matriz companheira; e a \textbf{membrana de Dirac} é a
raiz quadrada da transição do Metrónomo: não existe no próprio andar
(o dual proíbe, $\det^2=-1$), existe \emph{inteira} um andar acima,
aterra na diagonal exatamente nas transições, e o seu espectro vive na
\textbf{borda dobrada} $x^{4}-mx^{2}-1$ --- a dobra temporal
$x\mapsto x^{2}$ a aceder representantes na dimensão acima. O
\textbf{Corpo de Peano} (\code{papers/corpo\_peano.tex}) é
\emph{instância} deste quadro: a sua realização dinâmica e musical
(Maestro, Metrónomo, Batuta, Orquestra, Conservatório). Os papéis
ganham cidadania \emph{sob condição}: cada nome é o papel funcional que
um medidor validou. Verificação:
\code{tests/dirac\_transicao.js} ($13$:$0$),
\code{tests/ponte\_universal\_metalica.js} ($12$:$0$),
\code{tests/observador\_torre.js} ($10$:$0$),
\code{tests/assinatura\_colisoes.js}, \code{tests/assinatura\_banal.js},
\code{tests/cristal\_energia.js}, \code{tests/cristal\_volta.js}.
\end{abstract}

\medskip
\noindent\textbf{Leitura.}
\[
\boxed{
\text{Corpo Universal (este paper)}
\;\Longrightarrow\;
\text{Corpo de Peano (instância)}
\;\Longrightarrow\;
\text{realização musical / tradutor}.
}
\]
Nada aqui funda as oito leis: o catálogo é o do \emph{Corpo estelar}, e
este paper organiza o que os medidores fecharam por cima dele.

%======================================================================
\section{As cinco operações --- cada uma com a sua conservação}\label{sec:cinco}

\begin{definicao}[as cinco operações do corpo universal]\label{def:cinco}
\begin{center}\small
\begin{tabular}{@{}lll@{}}
\toprule
operação & conservação medida & medidor \\
\midrule
\textbf{Soma} $\oplus$ (a dobra/clone, Lei~4) & $E(u\oplus u)=2E(u)$; a retração devolve & \code{observador\_torre} \S T5 \\
\textbf{Multiplicação} $\otimes$ (fusão) & $E(u\otimes v)=E(u)E(v)$ --- $N(xy)=N(x)N(y)$ & \code{observador\_torre} \S T3 \\
\textbf{Divisão} (a fibra) & $x_i=z_{ij}/v_j$ exata; $E(u)=E(z)/E(v)$ & \code{observador\_torre} \S T4 \\
\textbf{Dual} $\dagger$ (a involução) & $\sigma\sigma^{\dagger}=-1$; $\lambda^{+}+\lambda^{-}=0$ & \code{lyapunov\_refletido}; \emph{Peano} \\
\textbf{Inversão} (a volta) & admissível $\iff$ tem volta $\wedge$ conserva a escada & \code{ponte\_universal\_metalica} \S P3 \\
\bottomrule
\end{tabular}
\end{center}
A Soma conserva \emph{adicionando}; a Multiplicação conserva
\emph{multiplicando} --- o par aditivo/multiplicativo da torre; a
Divisão é a \textbf{fibra inversa} da fusão (dividir custa exatamente o
$\det$: $\operatorname{adj}(M)\,M=\det(M)\,I$); o Dual guarda a segunda
metade; a Inversão é o critério de admissibilidade.
\end{definicao}

\begin{definicao}[o Inversor]\label{def:inversor}
Uma transformação $T$ é \textbf{admissível} quando
\[
\boxed{
\exists\,T^{-1}\ \text{(a volta)}
\qquad\land\qquad
I\bigl(T^{-1}(T(x))\bigr)=I(x)\ \text{na escada completa}.
}
\]
O \textbf{Inversor} é o papel funcional que decide: mede a volta contra
$(E,\Phi,\Phi_2,\ldots)$ e rejeita no primeiro desvio. Medido: a
transformada metálica passa (volta inteira nas unidades, escada
idêntica); a membrana de Dirac passa
(\code{tests/dirac\_transicao.js} \S D5).
\end{definicao}

%======================================================================
\section{O observador --- a escada, citada da instância}\label{sec:obs-uni}

A escada $I_1=(E)\prec I_2=(E,\Phi)\prec I_3=(E,\Phi,\Phi_2)\prec\cdots$
e os seus teoremas --- a estrita é orgânica ($4286/4286$ colisões da
dupla transposição espelhada; $\Phi_2$ separa todas), o \textbf{teto
pelas oito leis} (oito momentos fecham suporte $\le8$; $\Delta^{8}$
escapa ao linear e cai na energia --- a cruz) --- estão enunciados e
medidos na instância (\emph{Corpo de Peano}, secções do observador e da
energia). O universal só fixa o papel:
\[
\boxed{
\text{cada observador define uma resolução de identidade;}
\quad
\text{o \textbf{Conservatório} é o espaço dos invariantes.}
}
\]
\emph{Conservar a norma não conserva a identidade} --- por isso o fecho
é sempre um \textbf{par} de conservações concorrentes
($R_{\mathrm{total}}=(R_{\mathrm{endereço}},R_E)=0$), e a escada mede a
ordem sem ser a definição dela.

%======================================================================
\section{A membrana de Dirac --- a raiz da transição}\label{sec:dirac}

A transformada universal (avaliação nas folhas $\sigma,\sigma^{\dagger}$;
realização inteira pela matriz companheira --- \emph{Corpo de Peano},
a ponte) tem agora a sua \textbf{construção de Dirac}: o operador que
fatora a transição do Metrónomo. A escolha do coordenador --- a raiz
quadrada, \emph{metade para cada lado} --- e a derivação que já existia:

\begin{teorema}[a borda é uma fatoração de Dirac]\label{thm:borda-dirac}
No próprio andar, $D=A_m$ satisfaz
\[
\boxed{D^{2}=m\,D+I}
\]
--- a equação da borda $x^{2}=mx+1$ \emph{é} a fatoração: o passo é a
raiz quadrada do salto composto $L=A_m^{2}$. E \emph{metade para cada
lado}: $A^{k}A^{k}=A^{2k}$ com $\det(A^{k})^{2}=\det(A^{2k})=+1$ ---
cada metade carrega $\det=(-1)^{k}$ e o par fecha
($\sigma\sigma^{\dagger}=-1$).
\emph{Medido:} \code{tests/dirac\_transicao.js} \S D1, $m=1..6$, exato.
\end{teorema}

\begin{teorema}[a obstrução, e o furo na torre]\label{thm:furo}
A raiz quadrada do \emph{próprio passo} não existe no andar:
\[
\boxed{
\nexists\,D\in M_2(\mathbb{Z})\ \text{com}\ D^{2}=A_m
}
\qquad
\text{(}\det(D)^{2}=-1\ \text{impossível; }\sigma^{\dagger}<0\text{)},
\]
--- a obstrução é o \textbf{preço do dual} $\sigma\sigma^{\dagger}=-1$.
E existe \textbf{inteira, um andar acima}: na dobra ($4\times4$),
\[
\boxed{
D=\begin{pmatrix}0&I\\A_m&0\end{pmatrix},
\qquad
D^{2}=A_m\oplus A_m,
\qquad
\det D=-1\ \text{(unidade: inversa inteira).}
}
\]
\emph{Medido:} busca exaustiva vazia no $2\times2$; $D^{2}=A\oplus A$
exato, $m=1..6$; $D^{-1}D=I_4$ (\S D2--D3).
\end{teorema}

\begin{teorema}[Dirac atua exatamente nas transições]\label{thm:dirac-transicao}
$D$ aterra na diagonal $\{(z,z)\}$ \textbf{se e só se} o par é uma
transição verdadeira do Metrónomo:
\[
\boxed{
D(x,y)\in\text{diagonal}
\iff
y=A_m x.
}
\]
As transições são o lugar próprio do operador --- e a diagonal é o seu
conjunto fixo de aterragem. O Inversor aceita $D$: a volta é inteira e
a escada $(E,\Phi,\Phi_2)$ regressa idêntica.
\emph{Medido:} $100/100$ transições na diagonal, $100/100$
não-transições fora; $30$ corpos com a escada restaurada (\S D4--D5).
\end{teorema}

\begin{teorema}[o salto espectral --- a dobra temporal]\label{thm:dobra-temporal}
$D$ satisfaz
\[
\boxed{D^{4}=m\,D^{2}+I}
\qquad\text{isto é}\qquad
\boxed{\text{mín}(D)=x^{4}-mx^{2}-1=\bigl(x^{2}\bigr)^{2}-m\bigl(x^{2}\bigr)-1:}
\]
o polinômio mínimo de $D$ é a \textbf{borda dobrada} pela substituição
$x\mapsto x^{2}$ --- a \textbf{dobra temporal}: dois meios-passos são um
tick. Os representantes espectrais de $D$ ($\pm\sqrt{\sigma}$,
$\pm\sqrt{\sigma^{\dagger}}$) vivem no corpo de grau $4$ --- \textbf{o
andar acima da torre}, acedido pela dobra. E a leitura adjunta fecha o
outro lado: $D^{\dagger}D=(mA+I)\oplus I$ --- o salto inteiro num lado,
a identidade no outro: \emph{metade para cada lado}, literalmente.
\emph{Medido:} \S D6, $m=1..6$, exato.
\end{teorema}

\begin{definicao}[membrana]\label{def:membrana}
Uma \textbf{membrana} é um corpo (qualquer função endereçada) que
atravessa a transformada. A sua \textbf{representação espectral no
corpo universal} obtém-se pelas \textbf{dobras temporais}: cada dobra
$x\mapsto x^{2}$ sobe um andar e a membrana ganha os representantes
desse andar (Teorema~\ref{thm:dobra-temporal}). Estatuto:
\emph{definição de realização} sobre os teoremas medidos --- a membrana
não é operador novo; é o corpo em trânsito pela fatoração.
\end{definicao}

%======================================================================
\section{Os papéis --- a nomenclatura com cidadania}\label{sec:papeis}

Os nomes deixaram de ser abstratos: cada um é o papel funcional que um
medidor validou.
\begin{center}\small
\begin{tabular}{@{}lll@{}}
\toprule
papel & função no universal & onde o medidor fechou \\
\midrule
\textbf{Metrónomo} & métrica/ordem: lê a transição & \S D0, \S D4; \emph{Peano} (Teor.\ do Metrónomo) \\
\textbf{Maestro} & dinâmica: projecta o passo & \emph{Peano} (Teor.\ do Maestro) \\
\textbf{Inversor} & a volta: admissibilidade & \S D5; \S P3 \\
\textbf{Orquestra} & composição dos operadores & \emph{Peano} (Estrela dos $G_i$) \\
\textbf{Conservatório} & o espaço dos invariantes & a escada $(E,\Phi,\Phi_2)$; \S B4 \\
\textbf{Membrana} & o corpo em trânsito & Def.~\ref{def:membrana} \\
\bottomrule
\end{tabular}
\end{center}

%======================================================================
\section{O Corpo de Peano é instância}\label{sec:instancia}

\[
\boxed{
\text{Corpo de Peano}\;\prec\;\text{Corpo Universal}:
}
\]
o Peano realiza o quadro com escolhas concretas --- a torre
$d_k=2^{k+1}$, o catálogo $\mathbb{Z}/8\mathbb{Z}$, a retração $\pi_k$,
o par Maestro/Metrónomo, a realização musical
(partitura$\to$orquestra), a histerese e a rede dual. O universal não
o repete: \textbf{cita-o} --- cada teorema de lá é um caso deste
quadro, e cada teorema daqui foi medido sobre os corpos de lá (o
cristal, a família metálica, o corpus banal).
\begin{center}\small
\begin{tabular}{@{}lll@{}}
\toprule
universal & instância no Peano & medidor \\
\midrule
Soma $\oplus$ & $T_{k+1}=T_k+T_k^{*}$ (Lei~4) & \code{observador\_torre} \S T5 \\
Multiplicação $\otimes$ & viveiro ($\operatorname{lcm}\cdot\operatorname{gcd}=ab$) & \S T3 \\
Divisão/fibra & retração $\pi_k$; $\operatorname{adj}\cdot M=\det\cdot I$ & \S T4; \S P2 \\
Dual & estaca $\nu$; $\sigma^{\dagger}$ & \emph{Peano}, Lei~1 \\
Inversão & Metrónomo ($\lambda^{+}+\lambda^{-}=0$) & \code{lyapunov\_refletido} \\
escada $I_k$ & \S observador do \emph{Peano} & \S B4, \S C1--C2 \\
membrana & o corpo que o tradutor compõe & \code{cristal\_volta} \\
\bottomrule
\end{tabular}
\end{center}
E as \textbf{oito leis} continuam sendo o catálogo de ambos: o teto do
observador é o teto das oito (oito momentos fecham suporte $\le8$), e o
fecho $\operatorname{Ind}^{8}=\mathrm{id}$ vale no universal porque vale
no catálogo --- não há lei nova em nenhum dos dois.

%======================================================================
\section{Primitivas operacionais mecânicas}\label{sec:primitivas}

As engrenagens do sistema, cada uma \textbf{biunívoca} com um medidor
verde --- a regra de ouro da mecânica:
\[
\boxed{
\text{o que tem código verde entra; o que for só linguagem fica fora.}
}
\]
\begin{center}\small
\begin{tabular}{@{}clll@{}}
\toprule
 & primitiva & o que faz & medidor (verde) \\
\midrule
1 & \textbf{Núcleo de leitura} & corpo \& endereço: extrai o cristal;
  a cardinalidade conta a carga & \code{cristal\_volta.js} $16$:$0$ \\
2 & \textbf{Resolução de ordem} & a escada $(\Phi,\Phi_2)$: impede a
  cegueira da energia isolada & \code{assinatura\_colisoes.js};
  \code{assinatura\_banal.js} \\
3 & \textbf{Transformação universal} & família metálica:
  $\det=\pm1$, inversa inteira exata & \code{ponte\_universal\_metalica.js}
  $12$:$0$ \\
4 & \textbf{Fibra de desagregação} & divisão: decompositores admissíveis
  por $F^{-1}(z)$ sob conservação & \code{observador\_torre.js} \S T4 \\
5 & \textbf{Membrana de Dirac} & transição entre estados; fatoração e
  subida de andar quando a raiz não cabe & \code{dirac\_transicao.js}
  $13$:$0$ \\
6 & \textbf{Inversor total} & $R_{\mathrm{total}}=(R_{\mathrm{end}},R_E,
  R_\Phi,R_{\Phi_2},R_D)$: RETAIN $\iff$ tudo zera & \code{tests/residuos\_totais.js}
  $9$:$0$ \\
\bottomrule
\end{tabular}
\end{center}
As seis compõem a máquina completa: ler pelo endereço, resolver a
ordem, transformar com volta, desagregar pela fibra, atravessar a
transição, e admitir só o que conserva tudo --- as \textbf{operações
mecânicas} de que as cinco operações algébricas
(\S\ref{sec:cinco}) são a leitura.

%======================================================================
\section{Os domínios como realizações das oito leis}\label{sec:dominios}

O inventário do repo, com o estatuto medido --- \textbf{a tabela não
afirma equivalência entre domínios}: afirma que cada domínio realiza
operações do catálogo com resíduo $0$, no seu próprio medidor. A ponte
\emph{entre} domínios continua a construir-se elo a elo.
\begin{center}\small
\begin{tabular}{@{}lllll@{}}
\toprule
domínio & objeto no repo & leis realizadas & medidor (atestado) & estado \\
\midrule
álgebra & família metálica $A_m$ & 4, 7 (mult./inversão) &
  \code{ponte\_universal\_metalica.js} & verde \\
análise & transformada dourada; Parseval/Lebesgue & 8 (anel); central &
  \code{cristal\_energia.js} & verde \\
topologia & a topologia dos corpos: $d(r_1,r_2)=|\Delta_1-\Delta_2|$;
  $d=0\iff$ isomorfos; transporte parabólico $\det 1$ & 1, 2 &
  \code{tests/topologia.c} & verde \\
geometria & o simplex é a \emph{sombra} da base do andar acima
  (projeção dual e reversível); cone/espiral & 4, 5 &
  \code{tests/sombra\_cone.c}; \code{cone\_espiral.c} & verde \\
termodinâmica & Carnot pelo $\oint$, exato em $\mathbb{Q}$:
  $\eta=1-T_f/T_q$ derivado; o irreversível \emph{cai} & 0, 4
  (reverte ou paga) & \code{tests/carnot.c} & verde \\
cosmologia & a expansão completa em p.u., inteira;
  $\rho a^{3(1+w)}=C$; $w$ em órbitas de quatro; $w=-1$ ponto fixo &
  2, 4 & \code{tests/cosmologia.c} & verde \\
redes neurais & perceptron $=$ estaca ($W_P=-I$, Lei~1); rotor
  $J^{2}=-I$ (Lei~2); Hopfield $=$ memória da volta;
  $\lambda^{+}+\lambda^{-}=0$ & 1, 2, 3, 4 &
  \code{tests/rede\_dual.js} $28$:$0$ & verde \\
\quad$\hookrightarrow$ Clifford & $Cl(1,1)$ sobre o corpo metálico:
  $Df=N(f)f^{-1}$ (a estaca; graduação $\mathbb{Z}_2$), $\{D,G\}=0$,
  $D^{2}=L$, $G^{2}=-L$; $v^{2}=(x^{2}-y^{2})L$ emerge da soma dual &
  1, 2, 7 & \code{tests/clifford\_dual.js} $14$:$0$ & verde \\
\bottomrule
\end{tabular}
\end{center}
\begin{observacao}[o estatuto da tabela]
Cada linha verde foi \textbf{re-atestada} na data deste texto (os
medidores C compilados e corridos; os js com \code{\#TOTAL} $0$
falhas). «Verde» diz: \emph{o domínio realiza as leis citadas no seu
medidor}. Não diz: «os domínios são a mesma coisa» --- essa afirmação
continua a exigir a sua ponte (\S\ref{sec:extensoes}). A frase do
sistema para a termodinâmica já estava escrita na Arquitetura:
\emph{todo passo reverte, ou paga} --- e o \code{carnot.c} mede
exatamente o pagar ($\oint dQ/T<0$ faz $\eta$ cair).
\end{observacao}

%======================================================================
\section{A morfologia --- o par $\delta\dashv\varepsilon$ e a torção}\label{sec:morfologia}

A morfologia do Universal não se importa de fora: \textbf{a construção
já estava no repo} (o mórfico do \emph{Catálogo}; \code{tests/toolkit.c}
\S K3; \code{tests/morfico\_ordem.c}; Porquês${}=\varepsilon$ /
5W2H${}=\delta$ no Conecthus), e o laboratório
(\code{tests/morfologia\_universal.js}, $14$:$0$) verificou-a pelo
critério decisivo:
\[
\boxed{
\text{morfologia válida}
\iff
\text{dual existe}\;\land\;\text{volta existe}\;\land\;\text{invariante conservado}.
}
\]

\begin{teorema}[o par morfológico]\label{thm:morf-par}
A dilatação e a erosão são o par \textbf{adjunto}
\[
\boxed{\delta(A)\subseteq C\iff A\subseteq\varepsilon(C)}
\qquad(4096/4096\ \text{pares, varredura completa}),
\]
e dele \textbf{emergem} --- não se postulam --- a abertura e o fecho
clássicos: $\alpha=\delta\circ\varepsilon\le\mathrm{id}\le
\varepsilon\circ\delta=\varphi$, ambos \emph{idempotentes}. A
\textbf{volta existe na fronteira admissível}:
$\delta(\varepsilon(A))=A$ \emph{exatamente} nos conjuntos
$B$-abertos ($29$ de $64$ no universo medido) --- $E(D(x))=x$ dentro
da fronteira, como a adjunção manda. O invariante é a \textbf{ordem do
parâmetro}: o raio \emph{soma}
($\operatorname{dil}_r\circ\operatorname{dil}_s=\operatorname{dil}_{r+s}$,
Minkowski), total e monótona no objeto. E o par \textbf{muda o corpo}:
$\delta$ cria ($A\subseteq\delta A$), $\varepsilon$ remove
($\varepsilon A\subseteq A$) --- a massa sobe e desce, medida.
\end{teorema}

\begin{teorema}[a torção]\label{thm:torcao}
A torção tem \textbf{duas realizações que concordam}, e nenhuma é
dilatação escondida:
\begin{enumerate}
\item \textbf{o esquilo} $J\colon(a,b)\mapsto(b,-a)$:
$J^{2}=-I$, $J^{4}=\mathrm{id}$, $|\det J|=1$, e a \textbf{norma
conserva exata} --- $J$ transporta a coordenada para o eixo dual sem
criar nem remover;
\item \textbf{a órbita do anel}: a torção $w$ ($3^{256}$ em
$\mathbb{Z}_{65537}$) fecha em $256$ e em nenhum divisor próprio ---
\emph{a base é a órbita da torção} (\emph{Catálogo}) --- e \textbf{o
dual do dual devolve o grupo}: ida pelos caracteres, volta pela órbita
inversa, resíduo $0$ --- o degrau \emph{discreto} de Pontryagin,
medido no anel.
\end{enumerate}
\[
\boxed{
\text{erosão/dilatação mudam o \textbf{corpo}};
\qquad
\text{a torção muda o \textbf{representante} sem o destruir.}
}
\]
E a decomposição da morfologia universal é
\[
\boxed{
\text{morfologia}=\underbrace{\delta/\varepsilon}_{\text{escala/representação}}
+\underbrace{J}_{\text{eixo/assinatura}},
}
\qquad
\text{corpo}\xrightarrow{\delta/\varepsilon}\text{representação}
\xrightarrow{J}\text{dual}\xrightarrow{\text{Inversor}}\text{corpo}.
\]
A torção participa do par de Clifford recém-fechado ($J^{2}=-I$ é o
rotor da Lei~2; $\mathrm{Cl}(1,1)$ dá a assinatura) --- a morfologia lê
a escala, a torção assina o eixo, o Inversor confere a volta, o
Metrónomo mede o custo.
\emph{Medido:} \code{tests/morfologia\_universal.js} \S M1--M6;
inventário re-atestado: \code{toolkit.c} ($8$ ok),
\code{morfico\_ordem.c} ($3$ ok), \code{universal.c} ($42$ ok),
\code{rede\_dual.js} ($28$:$0$).
\end{teorema}

%======================================================================
\section{Extensões nomeadas --- o mapa, sem teorema}\label{sec:extensoes}

Cada nome abaixo \textbf{espera a sua ponte}; nenhum entra como teorema
antes do seu medidor:
\begin{itemize}
\item \textbf{Pontryagin} --- o eixo álgebra/topologia da teoria (a
dualidade dos caracteres); o degrau \emph{discreto} já está medido (o
bidual do anel devolve o grupo --- Teorema~\ref{thm:torcao}); o
\emph{contínuo} (o limite da escada de observadores) continua no mapa.
\item \textbf{Clifford} --- $D^{2}=L$ é a assinatura de uma relação de
Clifford; a estrutura medida é \emph{necessária} mas não suficiente
para a álgebra plena --- a próxima ponte natural.
\item \textbf{La Hire} --- o operador geométrico do rolamento dual (o
círculo de raio $r$ dentro de $2r$ fecha rotação em translação):
candidato geométrico do rotor da Lei~2.
\item \textbf{Dirac contínuo e convolução universal} --- registrados
como \emph{ausentes} da ponte metálica (\S P4 da ponte); a membrana
discreta (\S\ref{sec:dirac}) é o primeiro degrau deles.
\end{itemize}
\[
\boxed{
\text{a visão unificada é hipótese testável por elos ---}
\quad
\text{não se decreta; constrói-se ponte a ponte.}
}
\]

\bigskip
\noindent\textbf{Verificação.}
\code{tests/dirac\_transicao.js} \S D0--D6 ($13$:$0$);
\code{tests/residuos\_totais.js} \S R1--R6 (o Inversor total);
\code{tests/cristal\_adversarial.js} (a caçada às malditas);
\code{tests/cristal\_front.js} (o chip X);
\code{tests/equivalencia\_universal.js} (o adaptador $\mathcal U[\sigma]\cong\mathcal P$);
\code{tests/ponte\_universal\_metalica.js} \S P0--P5 ($12$:$0$);
\code{tests/observador\_torre.js} \S T1--T5 ($10$:$0$);
\code{tests/assinatura\_colisoes.js} \S C0--C4;
\code{tests/assinatura\_banal.js} \S B0--B5;
\code{tests/cristal\_energia.js} \S E0--E5;
\code{tests/cristal\_volta.js} \S V0--V4;
instância: \code{papers/corpo\_peano.tex} (§§ absorção, energia,
observador, ponte).

\end{document}
